برای این‌که نشان دهیم زبان $A$ یک زبان $\text{NP-complete}$ است، باید دو شرط زیر را اثبات کنیم:
\begin{enumerate}
    \item $A \in NP$
    \item برای هر زبان $L \in NP$، کاهش زمان چندجمله‌ای $L \le_P A$ وجود دارد.
\end{enumerate}

طبق فرض مسئله $A \in P$ است. می‌دانیم که همیشه $P \subseteq NP$ برقرار است. بنابراین به طور مستقیم داریم $A \in NP$.

فرض کنید $L$ یک زبان دلخواه در $NP$ باشد. طبق فرض $P = NP$، نتیجه می‌شود که $L \in P$ است. بنابراین، یک ماشین تورینگ قطعی مانند $M_L$ وجود دارد که زبان $L$ را در زمان چندجمله‌ای تصمیم‌گیری می‌کند.

از طرفی چون $A \neq \emptyset$ و $A \neq \Sigma^*$ است، حداقل یک رشته $w_{in} \in A$ و حداقل یک رشته $w_{out} \notin A$ وجود دارند.

حال تابع کاهش $f: \Sigma^* \to \Sigma^*$ را برای ورودی $x$ به صورت زیر تعریف می‌کنیم:
\[
f(x) = \begin{cases} 
w_{in} & x \in L  \\
w_{out} & x \notin L 
\end{cases}
\]

    بر اساس تعریف تابع $f$ داریم:
    \[
    x \in L \iff f(x) = w_{in} \in A
    \]
    \[
    x \notin L \iff f(x) = w_{out} \notin A
    \]
    بنابراین $x \in L \iff f(x) \in A$ همواره برقرار است.

    برای محاسبه $f(x)$، ابتدا ماشین $M_L$ را روی $x$ اجرا می‌کنیم که این کار طبق فرض $L \in P$ در زمان چندجمله‌ای انجام می‌شود. سپس بسته به نتیجه، یکی از دو رشته ثابت $w_{in}$ یا $w_{out}$ را در زمان $O(1)$ خروجی می‌دهیم. بنابراین تابع $f$ در زمان چندجمله‌ای قابل محاسبه است.

نتیجه می‌گیریم که $L \le_P A$ است و از آنجا که $L$ یک زبان دلخواه در $NP$ بود، شرط دوم نیز برقرار است.
